\documentclass[11pt]{article}
\usepackage{graphicx} % Required for inserting images
\usepackage{enumitem}
\usepackage{xcolor}
\usepackage{hyperref}
\usepackage{subcaption}
\usepackage{comment}
\usepackage{amsmath}
\usepackage{amssymb}
\usepackage{bm}
\usepackage{float}
\usepackage{listings}
\usepackage{multirow}
\usepackage{booktabs}
\usepackage[capitalise]{cleveref}
\usepackage{tikz}
\usetikzlibrary{arrows.meta, patterns}

\usepackage[margin=1in]{geometry}
\emergencystretch=2em % Lets long \texttt names wrap instead of overflowing the margin

% Items still missing before submission are highlighted in red
\newcommand{\missing}[1]{\textcolor{red}{[MISSING: #1]}}

% Includes the figure if the file exists; otherwise draws a red placeholder box with the expected path
\newcommand{\fig}[2]{%
    \IfFileExists{#2}%
        {\includegraphics[width=#1]{#2}}%
        {\fcolorbox{red}{white}{\begin{minipage}[c][3cm][c]{\dimexpr #1-2\fboxsep-2\fboxrule\relax}\centering\color{red}\small MISSING FIGURE\\{\scriptsize\texttt{\detokenize{#2}}}\end{minipage}}}%
}

% Placeholder box for a figure that has not been taken yet: \missingfig{width}{description}
\newcommand{\missingfig}[2]{%
    \fcolorbox{red}{white}{\begin{minipage}[c][3cm][c]{\dimexpr #1-2\fboxsep-2\fboxrule\relax}\centering\color{red}\small MISSING: figure of #2\end{minipage}}%
}

\lstset{
    language=Matlab,
    basicstyle=\ttfamily\footnotesize,
    keywordstyle=\color{blue},
    commentstyle=\color{green!50!black},
    stringstyle=\color{purple},
    numbers=left,
    numberstyle=\tiny\color{gray},
    frame=single,
    breaklines=true,
    columns=fullflexible,
    keepspaces=true,
    literate={°}{{$^\circ$}}1, % listings cannot read UTF-8 characters directly
}

\title{EECE7324: Laboratory Assignment 4 (Module 5)}
\author{Andrew Benincasa and Pau Sabater}
\date{November 2026}

% Table rows generated by our scripts (\input would add an empty row after the last \hline)
\makeatletter\newcommand{\inputtable}[1]{\@@input #1 }\makeatother

\begin{document}

\maketitle

%=====================================================================
\section*{Part A: M-PSK}
%=====================================================================
\subsection*{Transmitter}
\begin{enumerate}
    \item We started from the new scripts provided with this module instead of reusing our Module 4 versions.

    \item Yes. The number of symbols is fixed, and the script takes \texttt{nDataSymbols*log2(M)} bits from \texttt{bits} (\cref{lst:tx-bits}), so a higher $M$ needs more bits. \texttt{bits} only has 2184, so with 300 data symbols it works up to $M = 128$ (2100 bits) and fails for $M = 256$ (2400 bits).

\lstinputlisting[
    caption={Number of symbols and bits in \texttt{transmitterSC.m}.},
    label={lst:tx-bits},
    rangeprefix=, rangesuffix=, includerangemarker=true,
    linerange={nPilotSymbols-60;,nDataSymbols-60*5;,bitsPilot-log2(M));,bitsData-log2(M));},
]{Scripts_Module5/transmitterSC.m}

    \item \texttt{pskmod()} turns the bits into M-PSK symbols~\cite{pskmod}. It is called as \texttt{pskmod(bits, M, 0, "gray", InputType="bit")}: the bits, the modulation order, the phase offset in radians, the symbol order, and \texttt{InputType="bit"}, which tells it to take the bits in groups of $\log_2 M$, one group per symbol. The output is a column of complex symbols on the unit circle, $s_m = e^{j(2\pi m/M + \phi)}$, one per group of bits~\cite{pskmod}.

    \item Yes, it uses \texttt{"gray"}, which is also the default~\cite{pskmod}. To turn it off we pass \texttt{"bin"} instead~\cite{pskmod} (Step 5), and the receiver has to do the same.

    \item
    \begin{enumerate}[label=\alph*)]
        \item We added the symbol order as a parameter \texttt{symOrder} of \texttt{transmitterSC.m} and as an optional fourth input of \texttt{symbolMapping.m}, \texttt{"gray"} by default so the header call still works, and put it in the file name (\cref{lst:symorder}). This gave the files \texttt{gray/$M$psk\_25SR\_gray.WAVEFORM} and \texttt{bin/$M$psk\_25SR\_bin.WAVEFORM} for $M = 2, 4, 8, 16, 32$.

\begin{lstlisting}[caption={Our changes to \texttt{transmitterSC.m} and \texttt{symbolMapping.m}.}, label={lst:symorder}]
% transmitterSC.m
M = 2;             % Modulation order.
modScheme = "psk"; % Modulation scheme ("psk" or "qam").
symOrder = "gray"; % Symbol order ("gray" or "bin").
...
symbolsPilot = symbolMapping(bitsPilot, M, modScheme, symOrder);
...
symbolsData = symbolMapping(bitsData, M, modScheme, symOrder);
...
scatterplot(symbolsData); % Generated constellation.
title(sprintf("%d-%s (%s)", M, upper(modScheme), symOrder));
...
waveformFormatConversion(sFrame, sprintf("%d%s_%gSR_%s", M, modScheme, Rs / 1e6, symOrder));

% symbolMapping.m
function symbols = symbolMapping(bits, M, modScheme, symOrder)
    if nargin < 4
        symOrder = "gray";
    end
    ...
        symbols = pskmod(bits, M, 0, symOrder, InputType="bit");
\end{lstlisting}

        \item \Cref{fig:tx-psk} shows them, made with our script \texttt{checks/plotTxConstellations.m}. Gray and binary look the same, since only the bits of each point change, so we wrote the bits next to the points.
        \begin{enumerate}[label=\roman*.]
            \item The points get closer together: they are $360^\circ/M$ apart, only $11.25^\circ$ for 32-PSK.
            \item Yes. With the same power, less noise or phase error is enough to move a symbol to its neighbor: for 32-PSK, $5.6^\circ$ is already a wrong symbol.
            \item More SNR (more power or a shorter distance), a lower $M$ when the channel is bad, or QAM, which also uses the amplitude and spreads the points over the whole I/Q plane (Part B).
        \end{enumerate}
    \end{enumerate}

    \begin{figure}[H]
        \centering
        \fig{\linewidth}{figures/tx_psk.pdf}
        \caption{Generated M-PSK data constellations with Gray (top) and binary (bottom) order, with the bits of each point.}
        \label{fig:tx-psk}
    \end{figure}

    \item As in Module 4, the SG was at 2.4 GHz and 10 dBm, with Dual Arb ON at 75 MHz. On the DSO (Channel 1, 50 $\Omega$, DC coupling) we used 4 $\mu$s/div to look at the signal: the screen shows 40 $\mu$s, more than two frames of 16.32 $\mu$s, with three short dips per frame at the boundaries between header, pilot and data (\cref{fig:dso-psk}). We checked the spectrum with an FFT at 2.4 GHz and saved $10^6$ samples at 20 GSa/s (50 $\mu$s, enough for two complete frames). The saved samples do not depend on the timebase, only on the sample rate and the number of points: our wired captures were saved at 10 ns/div and still have 50 $\mu$s at 20 GSa/s, which our script \texttt{checks/capturePower.m} checks for every file. All the parameters of Steps 10 to 14 are computed from these files, not from the screen.
    \begin{enumerate}[label=\alph*)]
        \item Wired: a coaxial cable from the SG RF Out to Channel 1, at 500 mV/div. The received power is about 8 dBm in the signal band (our script \texttt{checks/capturePower.m}), so we do not expect errors.
        \item Wireless: one antenna on the SG and one on Channel 1, about \missing{distance, $\approx 1.1$ m} apart and pointing the same way (\cref{fig:setup-wireless}), because in Module 4 the errors started at about 114 cm. The DSO was at \missing{V/div}, and the received power was about \missing{dBm}.
    \end{enumerate}

    \begin{figure}[H]
        \centering
        \missingfig{0.7\linewidth}{the wireless setup at $\approx 1.1$ m, with the tape measure}
        \caption{Wireless setup: SG with one antenna on its RF Out, and the second antenna connected to Channel 1 of the DSO.}
        \label{fig:setup-wireless}
    \end{figure}

    \begin{figure}[H]
        \centering
        \begin{subfigure}{0.48\linewidth}\centering\missingfig{\linewidth}{wired 16-PSK Gray at 4 $\mu$s/div}\caption{Wired (500 mV/div).}\end{subfigure}\hfill
        \begin{subfigure}{0.48\linewidth}\centering\missingfig{\linewidth}{wireless 16-PSK Gray at 4 $\mu$s/div}\caption{Wireless.}\end{subfigure}
        \caption{DSO captures of 16-PSK with Gray coding at 4 $\mu$s/div: signal (yellow) and FFT (purple).}
        \label{fig:dso-psk}
    \end{figure}

\end{enumerate}

\subsection*{Receiver}
\begin{enumerate}[resume]
    \item We set \texttt{M}, \texttt{modScheme} and \texttt{symOrder} to the values of each capture and passed \texttt{symOrder} on to \texttt{symbolMapping()} and \texttt{symbolDemapping()} (\cref{lst:rx}). The rest already matched: 25 MBd, 60 pilot and 300 data symbols, and the same bits and header. The sampling frequency is different on purpose: 20 GSa/s, the one of the DSO (800 samples per symbol). Our script \texttt{checks/runReceiverAll.m} checks that all the captures are at 20 GSa/s and decode with no errors.

\begin{lstlisting}[caption={Our changes to \texttt{receiverSC.m}.}, label={lst:rx}]
M = 2;             % Modulation order.
modScheme = "psk"; % Modulation scheme ("psk" or "qam").
symOrder = "gray"; % Symbol order ("gray" or "bin"), same as in the transmitter.
...
filename = "captures/wired/6a_2psk_gray.mat"; % Update with your file name.
...
symbolsPilot = symbolMapping(bitsPilot, M, modScheme, symOrder);
...
symbolsData = symbolMapping(bitsData, M, modScheme, symOrder);
...
bitsReceived = symbolDemapping(symbolsDataReceived, M, modScheme, symOrder);
\end{lstlisting}

    \item It turns the received baseband signal back into one complex value per symbol: for each symbol, it takes the part of the signal where its pulse is and computes the inner product with the same SRRC pulse of the transmitter (\texttt{symbols(iSymbol) = pulse * ySymbol}). The inputs are the received baseband samples of one block \texttt{r}, the samples per symbol \texttt{sps} ($L = 800$), and the number of symbols \texttt{nSymbols} ($N$). The output is a column of $N$ complex symbols. The pulse, \texttt{rcosdesign(0.75, 10, sps)}, has $10L + 1 = 8001$ samples and unit energy~\cite{rcosdesign}. In mathematical notation, with $p[n]$ the pulse and $N_p = 10L+1$,
    \begin{equation}
        \hat{s}_i = \sum_{n=1}^{N_p} p[n]\, r\big[n + (i-1)L\big], \qquad i = 1, \dots, N,
    \end{equation}
    the discrete version of $\int r(t)\, p\big(t - (i-1)T_s\big)\, dt$. With no channel, it gives back the transmitted symbols within 0.5\% (our script \texttt{checks/correlatorLoopback.m}).

    \item \texttt{pskdemod()} uses Gray by default~\cite{pskdemod}, so it only matched our Gray captures. We added \texttt{symOrder} as an optional input, as in the transmitter (\cref{lst:demap}).

    \texttt{pskdemod()} does the opposite of \texttt{pskmod()}: for each received symbol it picks the closest of the $M$ points and returns its $\log_2 M$ bits~\cite{pskdemod}. For PSK the closest point is the one with the closest phase, so each point owns a slice of $360^\circ/M$. Our script \texttt{checks/demodNearest.m} checks that it always picks the closest point.

\begin{lstlisting}[caption={Our changes to \texttt{symbolDemapping.m}.}, label={lst:demap}]
function bits = symbolDemapping(symbols, M, modScheme, symOrder)
    if nargin < 4
        symOrder = "gray";
    end
    ...
        bits = pskdemod(symbols, M, 0, symOrder, OutputType="bit");
\end{lstlisting}

    \item \textbf{SNR.} As in Module 4, from the spectrum, but of the whole 50 $\mu$s capture (\cref{fig:psd}). The power in the signal band, $B = (1+\beta)R_s = 43.75$ MHz, is $P_S + P_N$. The noise floor $N_0$ is the median of the PSD 100 to 400 MHz away from the carrier; closer, there are the shoulders of the signal and an interferer. Then $P_N = N_0 B$ and $\text{SNR} = P_S/P_N$.

    \textbf{$\bm{E_b/N_0}$.} $P_S = E_s/T_s = E_b \log_2 M \cdot R_s$ and $P_N = N_0 B$, so
    \begin{equation}
        \text{SNR} = \frac{E_b \log_2 M \cdot R_s}{N_0 B} \quad \Longrightarrow \quad \frac{E_b}{N_0} = \text{SNR} \cdot \frac{B}{R_s \log_2 M} = \text{SNR} \cdot \frac{1.75}{\log_2 M},
    \end{equation}
    i.e., $E_b/N_0\,[\text{dB}] = \text{SNR}\,[\text{dB}] + 2.43 - 10\log_{10}(\log_2 M)$.

    \textbf{SER and BER} come from \texttt{receiverSC.m} (one frame: 300 symbols).

    Our script \texttt{checks/pskTable.m} computes all of it (\cref{tab:psk-gray,tab:psk-bin}). The wired SNR is 54 dB for every $M$, with no errors. \missing{Wireless columns and comment: rerun \texttt{checks/pskTable.m} with the 1.1 m captures. We expect a lower SNR than wired and some errors for 16- and 32-PSK.}

    \begin{table}[H]
        \centering
        \small
        \setlength{\tabcolsep}{4pt}
        \begin{tabular}{c|cc|cc|cc|cc}
            \hline
            \multirow{2}{*}{\shortstack{Modulation\\order}} & \multicolumn{2}{c|}{SNR (dB)} & \multicolumn{2}{c|}{$E_b/N_0$ (dB)} & \multicolumn{2}{c|}{SER} & \multicolumn{2}{c}{BER} \\
             & Wired & Wireless & Wired & Wireless & Wired & Wireless & Wired & Wireless \\
            \hline
            \inputtable{tables/psk_gray.tex}
        \end{tabular}
        \caption{M-PSK with Gray coded signals (one frame per capture: 300 symbols, $300\log_2 M$ bits).}
        \label{tab:psk-gray}
    \end{table}

    \begin{table}[H]
        \centering
        \small
        \setlength{\tabcolsep}{4pt}
        \begin{tabular}{c|cc|cc|cc|cc}
            \hline
            \multirow{2}{*}{\shortstack{Modulation\\order}} & \multicolumn{2}{c|}{SNR (dB)} & \multicolumn{2}{c|}{$E_b/N_0$ (dB)} & \multicolumn{2}{c|}{SER} & \multicolumn{2}{c}{BER} \\
             & Wired & Wireless & Wired & Wireless & Wired & Wireless & Wired & Wireless \\
            \hline
            \inputtable{tables/psk_bin.tex}
        \end{tabular}
        \caption{M-PSK without Gray coded signals (one frame per capture: 300 symbols, $300\log_2 M$ bits).}
        \label{tab:psk-bin}
    \end{table}

    \begin{figure}[H]
        \centering
        \fig{\linewidth}{figures/psd_16psk_gray.pdf}
        \caption{PSD of the wired and wireless 16-PSK Gray captures, with the signal band (green) and the band used for the noise floor (gray). Made with our script \texttt{checks/plotSpectrum.m}.}
        \label{fig:psd}
    \end{figure}

    \item No. Each value comes from one frame (300 symbols), and with 0 errors we can only say, with 95\% confidence, that $\text{SER} < 3/300 = 1\%$ (``rule of three''). The same bound holds for the BER, since each symbol error is at least one bit error and the bits of a symbol fail together. To make it meaningful we would need:
    \begin{itemize}
        \item Many more symbols, about 100 errors' worth: many captures (one cannot be longer than 50 $\mu$s) and a longer random bit sequence.
        \item A lower SNR, where errors happen (larger distance or lower SG power), swept to get the BER vs.\ $E_b/N_0$ curve.
        \item Repetitions at different times and positions.
    \end{itemize}
    \missing{Update with the 1.1 m captures: number of errors and their bounds.}

    \item A bit error happens when a received symbol crosses into the slice of a neighbor, so what matters is how far the symbols are from the slice boundaries (\cref{fig:rx-psk}, made with our script \texttt{checks/plotRxConstellations.m}):
    \begin{itemize}
        \item \textbf{Noise and interference} turn each point into a cloud; over the air the clouds are clearly bigger than with the cable.
        \item \textbf{Phase errors} rotate the points along the circle, which is what matters for PSK. The receiver corrects a static phase from the header; a frequency offset, like the one in Module 4, would rotate the whole constellation.
        \item \textbf{Multipath (ISI)} without an equalizer also spreads the clouds.
        \item \textbf{The modulation order}: the same cloud is closer to the boundaries when the slices are narrower. Over the air, it fills 2\% of the half-slice of 2-PSK but 38\% of the one of 32-PSK.
        \item \textbf{The mapping}: with Gray coding an error to a neighbor costs only one bit.
    \end{itemize}
    Amplitude errors move the points across the circle and do not change the decision for PSK. \missing{Wireless: update with the 1.1 m captures.}

    \begin{figure}[H]
        \centering
        \begin{subfigure}{\linewidth}\centering\fig{\linewidth}{figures/rx_psk_wired.pdf}\caption{Wired.}\end{subfigure}

        \vspace{0.5em}
        \begin{subfigure}{\linewidth}\centering\fig{\linewidth}{figures/rx_psk_wireless.pdf}\caption{Wireless.}\end{subfigure}
        \caption{Received M-PSK data constellations with Gray coding (blue) and the ideal points (red). The received symbols of each capture are only divided by their mean amplitude, so that they can be drawn on top of the ideal points.}
        \label{fig:rx-psk}
    \end{figure}

    \item We cannot compare them with our captures: the SER and BER are 0 in both cases (\cref{tab:psk-gray,tab:psk-bin}). We expect the same SER, because the mapping does not move the points or the decision regions, and a lower BER with Gray coding, because most errors go to a neighbor, which with Gray coding is only one bit wrong ($\text{BER} \approx \text{SER}/\log_2 M$), while in binary order it can be several (\texttt{000} and \texttt{111} are neighbors in 8-PSK, \cref{fig:tx-psk}). \missing{Measured SER and BER with and without Gray coding from the 1.1 m captures (or a lower SG power), to check this.}

    \item \textit{Optional.} The EVM says how far the received symbols are from where they should be. For each symbol, the error vector goes from the ideal point to the received one, and the RMS EVM is the RMS length of the error vectors divided by the RMS amplitude of the ideal points, usually given in \% or dB~\cite{commevm}:
    \begin{equation}
        \text{EVM} = \sqrt{\frac{\sum_k |\hat{s}_k - s_k|^2}{\sum_k |s_k|^2}}.
    \end{equation}
    Unlike the SER and BER, it does not need errors to happen, and it includes everything that moves the symbols: noise, phase errors, ISI and distortion.

    We computed it with our script \texttt{checks/evmTable.m}, after multiplying the received symbols by one real factor so that they have the same average power as the transmitted ones (\cref{tab:evm}); it gives the same values as \texttt{comm.EVM}~\cite{commevm}. With the cable the EVM is about 2\% ($-34$ dB), and over the air 4 to 8\%. If there were only noise, the EVM would be $1/\sqrt{\text{SNR}}$, i.e., 0.2\% at the 54 dB of Step 10, so most of the 2\% must come from something other than that noise, probably the distortion that shows up as shoulders in \cref{fig:psd}. \missing{Wireless: update with the 1.1 m captures.}

    \begin{table}[H]
        \centering
        \small
        \begin{tabular}{c|cc|cc}
            \hline
            \multirow{2}{*}{$M$ (Gray)} & \multicolumn{2}{c|}{Wired} & \multicolumn{2}{c}{Wireless} \\
             & EVM (\%) & EVM (dB) & EVM (\%) & EVM (dB) \\
            \hline
            \inputtable{tables/evm.tex}
        \end{tabular}
        \caption{RMS EVM of the received data symbols (one frame, 300 symbols).}
        \label{tab:evm}
    \end{table}
\end{enumerate}

%=====================================================================
\section*{Part B: M-QAM}
%=====================================================================
\subsection*{Transmitter}
\begin{enumerate}[resume]
    \item
    \begin{enumerate}[label=\alph*)]
        \item We added the QAM branch with \texttt{qammod()}, with the same inputs as \texttt{pskmod()}: the bits, $M$, the symbol order and \texttt{InputType="bit"} (\cref{lst:qam})~\cite{qammod}. By default \texttt{qammod()} puts the points 2 apart, so the average power grows with $M$ (2, 6, 10 and 20 for $M = 4$, 8, 16 and 32). \texttt{UnitAveragePower=true} scales it to an average power of 1~\cite{qammod}, the same as \texttt{pskmod()} (our script \texttt{checks/qamPower.m}).
\begin{lstlisting}[caption={Our change to \texttt{symbolMapping.m} for M-QAM.}, label={lst:qam}]
    elseif strcmp(modScheme, "qam")
        symbols = qammod(bits, M, symOrder, InputType="bit", UnitAveragePower=true); % Average power 1, as pskmod.
\end{lstlisting}
    \end{enumerate}

    \item \Cref{fig:tx-qam}, made with our script \texttt{checks/plotTxQam.m} from \texttt{transmitterSC.m}, shows what we expected: 4-QAM is 4-PSK turned by $45^\circ$, 8-QAM is a $4 \times 2$ rectangle, 16-QAM a square and 32-QAM a cross, all with average power 1. With Gray coding, neighbors differ in one bit, except for 8 of the 52 pairs of neighbors in the 32-QAM cross. For $M \geq 8$ the points are farther apart than in M-PSK: the minimum distance is 0.82, 0.63 and 0.45 for $M = 8$, 16 and 32, against 0.77, 0.39 and 0.20.

    \begin{figure}[H]
        \centering
        \fig{\linewidth}{figures/tx_qam.pdf}
        \caption{Generated M-QAM data constellations with Gray coding, with the bits of each point.}
        \label{fig:tx-qam}
    \end{figure}

    \item
    \begin{enumerate}[label=\alph*)]
        \item With \texttt{modScheme = "qam"} and \texttt{symOrder = "gray"} in \texttt{transmitterSC.m}, we generated \texttt{qam/$M$qam\_25SR\_gray.WAVEFORM} for $M = 4$, 8, 16 and 32 (our script \texttt{checks/plotTxQam.m} runs it for the four). The frame is the same as for M-PSK: BPSK header, 60 pilot and 300 data symbols.
    \end{enumerate}

    \item We repeated Step 6 with the four M-QAM files, with the same SG settings and DSO setup.
    \begin{enumerate}[label=\alph*)]
        \item Wired: the captures are \texttt{wired/18a\_$M$qam\_gray.mat}. We did not write down the vertical scale; from the voltage step of the samples, our script \texttt{checks/capturePower.m} estimates about 420 mV/div. The power in the signal band is 8.1 dBm, the same as for M-PSK, but the peaks are higher (about 3.0 V peak to peak against 2.5 V), because QAM uses several amplitudes and the SG keeps the same average power. \missing{Screenshot of wired 16-QAM at 4 $\mu$s/div.}
        \item Wireless: \missing{same setup as Step 6b at $\approx 1.1$ m: files \texttt{wireless/18b\_$M$qam\_gray.mat}, V/div, received power, and a screenshot of 16-QAM.}
    \end{enumerate}
\end{enumerate}

\subsection*{Receiver}
\begin{enumerate}[resume]
    \item
    \begin{enumerate}[label=\alph*)]
        \item We added the QAM branch with \texttt{qamdemod()}, with the same symbol order and \texttt{UnitAveragePower=true} as in the transmitter, and \texttt{OutputType="bit"} (\cref{lst:qamdemod})~\cite{qamdemod}. Unlike for PSK, the QAM decision depends on the amplitude, and the received symbols have whatever amplitude the channel and the DSO give them, so we first divide them by their RMS value to give them average power 1, like the constellation. It is one real factor for all the symbols, so I and Q keep their ratio. Our script \texttt{checks/runReceiverAll.m} runs \texttt{receiverSC.m} on the 8 QAM captures, with no errors.
\begin{lstlisting}[caption={Our change to \texttt{symbolDemapping.m} for M-QAM.}, label={lst:qamdemod}]
    elseif strcmp(modScheme, "qam")
        % QAM decisions depend on the amplitude, so we first give the received symbols unit average
        % power, like the constellation (one real factor; no equalization until Module 6).
        symbols = symbols / sqrt(mean(abs(symbols).^2));
        bits = qamdemod(symbols, M, symOrder, OutputType="bit", UnitAveragePower=true);
\end{lstlisting}
    \end{enumerate}

    \item \Cref{fig:rx-qam} shows the received symbols as they come out of the receiver, without normalizing I or Q, made with our script \texttt{checks/plotRxQam.m}. The red crosses are the ideal points scaled by the RMS value of the received symbols. With the cable, the constellations have the shapes of \cref{fig:tx-qam}, with tight clusters (EVM of 1 to 2\%) and no rotation, so the phase correction from the BPSK header also works for QAM. Over the air, the points are 34 times smaller (31 dB, like the received power) and the clusters larger (EVM of 8 to 11\%), but they stay inside their regions, which is why there are no errors (Step 19). In the wired 16-QAM, the outer points are slightly outside the red crosses: the RMS scaling is a bit off because our 300 data symbols have an average power of 0.94 instead of 1 (Step 15). \missing{Wireless: update with the 1.1 m captures.}

    \begin{figure}[H]
        \centering
        \begin{subfigure}{\linewidth}\centering\fig{\linewidth}{figures/rx_qam_wired.pdf}\caption{Wired.}\end{subfigure}

        \vspace{0.5em}
        \begin{subfigure}{\linewidth}\centering\fig{\linewidth}{figures/rx_qam_wireless.pdf}\caption{Wireless.}\end{subfigure}
        \caption{Received M-QAM data constellations (blue), without normalizing I or Q, and the ideal points scaled by the RMS value of the received symbols (red).}
        \label{fig:rx-qam}
    \end{figure}

    \item We used the same method as in Step 10 (SNR from the spectrum, the same $E_b/N_0$ expression, and \texttt{receiverSC.m} with the QAM demapping of Step 19), with our script \texttt{checks/pskTable.m} (\cref{tab:qam-gray}). With the cable the SNR is 55 dB, a bit higher than for M-PSK because the DSO was at a smaller vertical scale (Step 18), and there are no errors. \missing{Wireless columns and comment: rerun \texttt{checks/pskTable.m} with the 1.1 m captures.}

    \begin{table}[H]
        \centering
        \small
        \setlength{\tabcolsep}{4pt}
        \begin{tabular}{c|cc|cc|cc|cc}
            \hline
            \multirow{2}{*}{\shortstack{Modulation\\order}} & \multicolumn{2}{c|}{SNR (dB)} & \multicolumn{2}{c|}{$E_b/N_0$ (dB)} & \multicolumn{2}{c|}{SER} & \multicolumn{2}{c}{BER} \\
             & Wired & Wireless & Wired & Wireless & Wired & Wireless & Wired & Wireless \\
            \hline
            \inputtable{tables/qam_gray.tex}
        \end{tabular}
        \caption{M-QAM with Gray coded signals (one frame per capture: 300 symbols, $300\log_2 M$ bits).}
        \label{tab:qam-gray}
    \end{table}

    \item With the same average power, M-QAM spreads its points over the whole I/Q plane, so for $M \geq 8$ they are farther apart than in M-PSK (Step 16), and we expect fewer errors for the same $E_b/N_0$: the minimum distance is larger by 0.6, 4.2 and 7.2 dB for $M = 8$, 16 and 32. For $M = 4$ both are the same constellation. The price is that QAM also depends on the amplitude, so it needs the right scaling (Step 19), and its peaks are higher (Step 18).

    We cannot see this in the SER and BER, since both are 0 (\cref{tab:psk-gray,tab:qam-gray}). The EVM shows it instead: what matters is the half distance to the neighbors compared with the error vectors, $d_\text{min}/(2\,\text{EVM})$. With the cable, for $M = 32$, the boundaries of 32-QAM are 16 RMS error vectors away (EVM of 1.4\%, \texttt{checks/plotRxQam.m}), against about 3 for 32-PSK (EVM of 3.4\%, \cref{tab:evm}). \missing{Wireless comparison with the 1.1 m captures: SER and BER of M-QAM and M-PSK at the same distance, and the same margin.}

    \item \textit{Optional.} The PAPR is the highest instantaneous power of a signal divided by its average power, usually in dB. It tells how much room the DAC and the amplifier need above the average: with a high PAPR the amplifier has to work further from saturation, so it is less efficient, or the peaks get clipped and distorted. All the M-PSK symbols have the same power, so their PAPR is 0 dB, while the outer points of M-QAM are above the average.

    \Cref{tab:papr}, from our script \texttt{checks/paprTable.m}, gives the PAPR of the symbols. It does not simply grow with $M$: it depends on how far the outer points are compared with the average, which is the most for 16-QAM, with its corners, and less for the 32-QAM cross, which has no corners. The received constellations have a slightly higher PAPR than the transmitted symbols, because the noise and distortion spread the clusters and the most extreme symbol sets the peak: by about 0.1 dB with the cable and 0.5 to 1.5 dB over the air. The PAPR of the waveform is higher still, since the pulse shaping adds peaks between the symbols, which is why the QAM captures have larger peaks than the M-PSK ones (Step 18). \missing{Wireless: update with the 1.1 m captures.}

    \begin{table}[H]
        \centering
        \small
        \begin{tabular}{c|cc|cc}
            \hline
            \multirow{2}{*}{$M$} & \multicolumn{2}{c|}{Transmitted (dB)} & \multicolumn{2}{c}{Received (dB)} \\
             & All $M$ points & 300 data symbols & Wired & Wireless \\
            \hline
            \inputtable{tables/papr.tex}
        \end{tabular}
        \caption{PAPR of the M-QAM symbols.}
        \label{tab:papr}
    \end{table}
\end{enumerate}

%=====================================================================
\section*{Part C: Matched Filter Receiver}
%=====================================================================
\begin{enumerate}[resume]
    \item We copied \texttt{functions/correlationReceiver.m} to \texttt{functions/matchedFilter.m} and renamed the function inside it to \texttt{matchedFilter}, since MATLAB expects the function to have the name of its file.

    \item We replaced only the line that correlates and integrates (\texttt{symbols(iSymbol) = pulse * ySymbol}): now the same piece of the received signal goes through the matched filter, the time-reversed pulse $h[n] = p[N_p - n + 1]$, with \texttt{conv()}, and we keep the output at the sample where the pulse and the filter overlap completely, which is the sampling instant $T_s$ of that symbol (\cref{lst:mf}). The option \texttt{"valid"} returns only that sample. In mathematical notation, with $r_i[n]$ the piece of the $i$-th symbol,
    \begin{equation}
        \hat{s}_i = (r_i * h)[N_p] = \sum_{n=1}^{N_p} r_i[n]\, h[N_p - n + 1] = \sum_{n=1}^{N_p} r_i[n]\, p[n].
    \end{equation}
\begin{lstlisting}[caption={Our change to \texttt{matchedFilter.m}, the only line that differs from \texttt{correlationReceiver.m} apart from the function name.}, label={lst:mf}]
% correlationReceiver.m
        symbols(iSymbol) = pulse * ySymbol; % Correlator receiver.
% matchedFilter.m
        symbols(iSymbol) = conv(ySymbol, fliplr(pulse).', "valid"); % Matched filter h[n] = p[Np - n + 1], sampled at Ts.
\end{lstlisting}

    \item They use the same filter. The SRRC pulse is symmetric (our script \texttt{checks/correlatorLoopback.m} checks that $p$ and its time-reversed version are identical), so the matched filter, the time-reversed pulse, is the pulse itself. What changes is the order: the transmitter upsamples the symbols and then filters them with $p$, going from symbols to samples, while the receiver filters with $p$ and then keeps one sample per symbol at $T_s$, going from samples to symbols. Together, the two SRRC filters give a raised cosine, which is zero at the other sampling instants, so there is no ISI (Step 8), and the matched filter gives the highest SNR at the sampling instants. The pulse has the same shape on both sides, but with 3 samples per symbol at the SG (75 MSa/s) and 800 at the DSO (20 GSa/s).
\end{enumerate}

\begin{thebibliography}{9}
    \bibitem{pskmod} MathWorks, ``pskmod: Modulate signal using M-PSK method,'' MATLAB Communications Toolbox documentation (R2026b). [Online]. Available: \url{https://www.mathworks.com/help/comm/ref/pskmod.html}
    \bibitem{rcosdesign} MathWorks, ``rcosdesign: Raised cosine FIR pulse-shaping filter design,'' MATLAB Signal Processing Toolbox documentation (R2026b). [Online]. Available: \url{https://www.mathworks.com/help/signal/ref/rcosdesign.html}
    \bibitem{pskdemod} MathWorks, ``pskdemod: Demodulate using M-ary PSK method,'' MATLAB Communications Toolbox documentation (R2026b). [Online]. Available: \url{https://www.mathworks.com/help/comm/ref/pskdemod.html}
    \bibitem{commevm} MathWorks, ``comm.EVM: Measure error vector magnitude (EVM) of received signal,'' MATLAB Communications Toolbox documentation (R2026b). [Online]. Available: \url{https://www.mathworks.com/help/comm/ref/comm.evm-system-object.html}
    \bibitem{qammod} MathWorks, ``qammod: Quadrature amplitude modulation (QAM),'' MATLAB Communications Toolbox documentation (R2026b). [Online]. Available: \url{https://www.mathworks.com/help/comm/ref/qammod.html}
    \bibitem{qamdemod} MathWorks, ``qamdemod: Quadrature amplitude demodulation,'' MATLAB Communications Toolbox documentation (R2026b). [Online]. Available: \url{https://www.mathworks.com/help/comm/ref/qamdemod.html}
\end{thebibliography}

\end{document}
